\item \textbf{¿Por qué no es posible la siguiente situación?}
Un fotón de longitud de onda $88.0\text{ nm}$ golpea una superficie de aluminio limpio, expulsando un fotoelectrón. Después, el fotoelectrón impacta a un átomo de hidrógeno en su estado fundamental, transfiriéndole energía al mismo y excitándolo a un estado cuántico más alto.